\section{Prefix maps with changing area}\label{sec:sa-sheets}
A stack operation changes a finite prefix and leaves the remaining word
unchanged.  In a positional code this becomes a diagonal affine map of
a rectangle.  The two scale factors need not be reciprocal.  We shall
therefore use either a normal strain to compensate the area change, or
a planar motion embedded in an invariant plane.  These two methods have
different geometric guarantees: the former moves a neighborhood in three
dimensions, whereas the latter preserves one plane throughout the motion.

Throughout these constructions, viscosity is a fixed positive computable
number.  At an arbitrary fixed positive real viscosity, the same formulas
are effective relative to that coefficient: the residual has form
$G+\nu J$ with effective coefficient fields independent of $\nu$.
Coordinate inequalities on the unit torus denote the displayed arcs modulo
one; the path estimates use the stated representatives inside a chart.

\subsection{From prefix cylinders to closed rectangles}\label{subsec:sa-local}
We use the machine convention of a two-sided tape, a finite alphabet with
blank, and moves $-1,0,1$.  A missing instruction interpreted as halting
can be replaced by an instruction that preserves the scanned symbol,
stays in place, and enters a new halt state.  Multiple halt states can
be merged when a construction calls for a single one.  A left-end
convention is represented by a permanent tag at the initial cell, with
writes preserving that tag and the prescribed endpoint rule used there.
These finite changes preserve the halting question.

A two-stack configuration is $(s,L,R)$, with both words written top first.
A rule $(s;v,w)\to(s';v',w')$ means
\[
 (s,vL_*,wR_*)\longmapsto(s',v'L_*,w'R_*)
\]
for arbitrary infinite tails.  A \emph{cylinder} is the set specified by
the state and the two finite prefixes.  We shall always check the full
source and image cylinders separately.  Agreement on initialized runs
alone would not justify a smooth invertible extension.

\begin{lemma}[Gapped stack coordinates]\label{lem:sa-radix}
Choose an integer $B\ge2$ and a finite digit set in
$\{0,\ldots,B-2\}$ whose distinct elements differ by at least two.
For words over the corresponding alphabet define
\[
 E(v)=\sum_{j=1}^{|v|}d(v_j)B^{-j},\qquad
 I(v)=[E(v),E(v)+B^{-|v|}],\qquad
 E(L)=\sum_{j\ge1}d(L_j)B^{-j}.
\]
The infinite code is injective and satisfies
$E(vL)=E(v)+B^{-|v|}E(L)$.  Incompatible finite prefixes give
positively separated closed intervals.  If separately disjoint state
rectangles are parameterized by positive diagonal affine maps $P_s$
from $[0,1]^2$, then separately disjoint source and image cylinders give
separately disjoint closed source and target rectangles
\[
 P_s(I(v)\times I(w)),\qquad P_{s'}(I(v')\times I(w')).
\]
The positive diagonal affine map between a matched pair implements the
rule on every encoded configuration in its source cylinder.
\end{lemma}
\begin{proof}
Prefixing a symbol selects the interval $(d+[0,1])/B$.  These intervals
lie in $[0,1]$, and two of them have a gap at least $B^{-1}$.
Further prefixing selects nested affine copies.  At the first difference
in position $j$, the gap is at least $B^{-j}$, which proves both interval
separation and injectivity.  The series identity follows by splitting
its first $|v|$ terms.  Two product cylinders with the same state are
disjoint precisely when the prefixes on at least one stack are
incompatible: otherwise their compatible prefixes have common infinite
extensions.  This proves rectangle separation.  Finally the affine map
preserves each normalized tail coordinate $E(L_*),E(R_*)$, which proves
its action on codes.  Codes of $v\square^\infty$ are rational, equal to
$E(v)+B^{-|v|}d(\square)/(B-1)$.  Digit gaps allow recovery of every
finite prefix to sufficiently high finite precision, including zero-digit
boundary codes.
\end{proof}

The history mechanism follows Bennett's information-retaining approach
\cite{Bennett1973}; the use of positional codes is closely related to
Moore's generalized shifts~\cite{Moore1990,Moore1991}.  The additional
work here is the full-cylinder inverse check and a smooth incompressible
extension with controlled intermediate trajectories.

\subsection{Local motion and normal compensation}
Write $\sigma$ for the effective flat step of \eqref{eq:p2-step} and
$\theta(s)=\sigma(3s-1)$.  All schedules below are finite rescalings of
these functions.  For an interval $[a,b]$ and margin $\rho>0$, the product
\[
 \sigma\!\left(\frac{x-a+\rho}{\rho/2}\right)
 \sigma\!\left(\frac{b+\rho-x}{\rho/2}\right)
\]
equals one on its $\rho/2$ enlargement and is supported in its $\rho$
enlargement.  Products in the coordinates give box cutoffs.  The effective
flatness estimates in Section~\ref{sec:model} justify evaluation of all
mixed derivatives across their seams; no equality test at a seam is used.

\begin{lemma}[Localized affine motions]\label{lem:sa-curl-motion}
Let $c(t)\in\mathbb R^3$ and a positive diagonal matrix $D(t)$ be smooth
effective functions on a compact interval, with $D(t_0)=I$ and
$\det D(t)=1$.  Put $A=\dot D D^{-1}$.  Let $K_0$ be a rectangular box, allowing
zero thickness in a coordinate.  If an effective smooth compactly supported
cutoff $\chi$ equals one on a
positive collar of the moving box $K_t=c(t)+D(t)(K_0-c(t_0))$, then
\begin{equation}\label{eq:sa-curl}
 V(t,x)=\nabla\times\left\{\chi(t,x)
 \left[\tfrac12\dot c\times(x-c)
       +\tfrac13(A(x-c))\times(x-c)\right]\right\}
\end{equation}
is divergence free and its trajectory from $x_0\in K_0$ is
\begin{equation}\label{eq:sa-affine-path}
 x(t)=c(t)+D(t)(x_0-c(t_0)).
\end{equation}
The formula holds on an effective positive neighborhood of $K_0$.
Disjoint moving supports allow finitely many such motions to be summed.
\end{lemma}
\begin{proof}
Differentiating $\det D=1$ gives $\operatorname{tr}A=0$.  With $Y=x-c$,
\[
 \nabla\times(\dot c\times Y)=2\dot c,\qquad
 \nabla\times((AY)\times Y)=3AY.
\]
The second identity follows by expanding the curl and using
$\operatorname{div}(AY)=0$.  Thus on the plateau $V=\dot c+AY$, and
\eqref{eq:sa-affine-path} solves its ODE.  If $M\ge1$ bounds $\|D\|$ and
the collar width is $\eta$, an initial perturbation smaller than
$\eta/(2M)$ remains on that plateau.  For a finite sequence, decrease
this radius by the product of the preceding affine norms.  Smooth ODE
uniqueness proves the assertion, and disjoint supports prevent interference.
\end{proof}

\begin{lemma}[Transport of area-changing sheets]\label{lem:sa-sheet-routing}
Let $B_i$ and $C_i$, $1\le i\le N$, be two separately positively
separated finite families of nondegenerate rational closed rectangles
in $(0,1)^2$.  Let $F_i:B_i\to C_i$ be the positive diagonal affine
map, and let $z_0\in(0,1)$ be rational.  There is an effective
smooth divergence-free field supported in $(0,1)\times(0,1)^3$ whose
time-one map agrees on a positive three-dimensional neighborhood of each
$B_i\times\{z_0\}$ with a diagonal affine map restricting to the
map $(x,z_0)\mapsto(F_i(x),z_0)$, with normal scale
reciprocal to its planar area multiplier.  On the
sheet, each horizontal coordinate is constant during ascent and descent
and is a convex combination of its two endpoint coordinates during the
middle stage.  The same construction applies in a larger bounded
Euclidean region after translations and a choice of separated heights.
\end{lemma}
\begin{proof}
For $N=0$ take zero.  Write $p_i,q_i$ for the centers and $w_i^0,w_i^1$
for side-length vectors.  In three successive stages, lift to
$z_i=1/2+i/[4(N+1)]$, transform there, and lower to $z_0$.  During the
middle stage, with a flat progress parameter $s\in[0,1]$, put
\[
 p_i(s)=(1-s)p_i+sq_i,\quad w_i(s)=(1-s)w_i^0+sw_i^1,
\]
\[
 D_i(s)=\diag\left(\frac{w_{i1}(s)}{w_{i1}^0},
 \frac{w_{i2}(s)}{w_{i2}^0},
 \frac{w_{i1}^0w_{i2}^0}{w_{i1}(s)w_{i2}(s)}\right).
\]
All scales are positive and their product is one.  Give the sheets a
sufficiently small rational thickness.  The reciprocal scale is bounded
by $\max_i[\min(1,w_{i1}^1/w_{i1}^0)
\min(1,w_{i2}^1/w_{i2}^0)]^{-1}$, so an effective positive thickness
keeps the middle-stage boxes apart vertically.  During lifting use the
source separation, and during lowering the target separation.  Choose a
rational collar smaller than a quarter of the corresponding coordinate
gaps and boundary clearances.  Lemma~\ref{lem:sa-curl-motion} now supplies
the three local motions on disjoint supports, including a positive initial
neighborhood.  A sheet point has zero normal displacement relative to its
center.  Its fixed fractional coordinate in the interpolated interval
gives $(1-s)x_a+sF_i(x)_a$ in each horizontal direction.  Flat stage
collars finish the construction.  The larger-region version uses exactly
the same finite clearance argument.
\end{proof}

\subsection{Closing the fluid argument}\label{subsec:sa-fluid-assembly}
\begin{lemma}[Realization and comparison]\label{lem:sa-realization}
Let $U$ be an effective smooth divergence-free field, $U(0)=0$, on the
unit torus or with fixed compact spatial support on $\mathbb R^3$.
Assume its mixed derivatives are bounded on each finite time interval.
Then
\begin{equation}\label{eq:sa-residual}
 f=\partial_tU+(U\cdot\nabla)U-\nu\Delta U
\end{equation}
is an effective prescribed force and $(u,p)=(U,0)$ is a global solution.
Its velocity is unique in the applicable comparison classes of
Lemma~\ref{lem:b-comparison}; pressure is determined with exactly the
normalization or additive time function specified there.
Uniform mixed-derivative bounds, a time period, fixed compact support,
and spatial mean zero pass from $U$ to $f$ when imposed on $U$.
The material flow is defined for every finite time.
\end{lemma}
\begin{proof}
The compactly supported field satisfies the whole-space velocity class
$C_tH^2\cap C_t^1L^2$ and has bounded velocity and gradient; zero pressure
satisfies $C_tH^1$.  Thus Lemma~\ref{lem:p2-realization} applies in either
domain.  Derivatives of \eqref{eq:sa-residual} are finite sums of products
of derivatives of $U$, proving the stated bounds and support.  For mean
zero use $(U\cdot\nabla)U=\operatorname{div}(U\otimes U)$ and integrate.
The effective cutoff formulas and differentiation give the force program.
Periodic evaluation uses a finite superset of potentially active translates,
including adjoining intervals at a seam.
\end{proof}

Two refinements of the whole-space comparison will be used explicitly.
Lemma~\ref{lem:b-comparison}(ii) permits pressure in $H^1$ modulo spatial
constants at each time, without continuity into that norm, when the
competing velocity retains $C_tH^2\cap C_t^1L^2$ and finite-interval
bounds on velocity and gradient.  Its case (i) instead permits omission
of the competing-gradient bound when pressure belongs to $C_tH^1$.
The compactly supported reference fields below satisfy the hypotheses
of both statements.  We retain each theorem's declared comparison class
when applying the corresponding case.

For each construction below, the finite rule list gives all spatial
branches before a trajectory is chosen.  The choice of initial coordinates
or a loading motion supplies the exact rational starting code; a loading
rule inside the repeated program is included in the induction at integer
times.  We then check the entire path of every used nonhalting branch
against its observer.  This last check is essential: a safe pair of
endpoints alone need not give a safe interpolation.  Bounds on the
finitely many pulse types give uniform kinetic energy on the torus,
or under fixed compact support in Euclidean space.  No simulation trace
enters a force prescription.
